\documentclass[aspectratio=169, 10pt]{beamer}
\usepackage{beamer_theme_cienciadados2026}

\title[Aula 09 - Pandas I]{Ciência dos Dados I: Análise Exploratória}
\subtitle{Aula 09: Análise de Dados com Pandas I: Series, DataFrames e Indexação}
\author{Prof. Dr. Ney Lemke}
\institute{UNESP -- Instituto de Biociências de Botucatu}
\date{Versão 2026}

\begin{document}

\frame{\titlepage}

\begin{frame}{Sumário da Aula}
  \tableofcontents
\end{frame}

\section{O Papel do Pandas no Ecossistema}

\begin{frame}{Por que Pandas?}
  Criada por Wes McKinney em 2008, o Pandas é a espinha dorsal da análise de dados tabulares em Python:
  \begin{itemize}
    \item Combina o desempenho de vetores contíguos do NumPy com a flexibilidade de planilhas e bancos relacionais.
    \item Suporta rótulos explícitos de linhas e colunas (índices).
    \item Facilita alinhamento automático de dados e tratamento de dados ausentes.
  \end{itemize}
  \begin{block}{Estruturas Centrais}
    \begin{itemize}
      \item \textbf{Series}: Matriz unidimensional rotulada capaz de conter qualquer tipo de dados.
      \item \textbf{DataFrame}: Estrutura de dados tabular bidimensional com eixos rotulados (linhas e colunas).
    \end{itemize}
  \end{block}
\end{frame}

\section{Indexação Rigorosa: \texttt{loc} vs. \texttt{iloc}}

\begin{frame}[fragile]{Seleção de Dados sem Ambiguidade}
  \begin{columns}
    \begin{column}{0.5\textwidth}
      \begin{block}{\texttt{.loc[linhas, colunas]}}
        \textbf{Baseada em RÓTULOS}:
        \begin{itemize}
          \item Usa nomes das colunas e índices.
          \item Intervalo inclui o ponto final!
        \end{itemize}
      \end{block}
      \begin{lstlisting}
df.loc[0:3, ['idade', 'glicose']]
      \end{lstlisting}
    \end{column}
    \begin{column}{0.48\textwidth}
      \begin{block}{\texttt{.iloc[linhas, colunas]}}
        \textbf{Baseada em POSIÇÃO INTEIRA}:
        \begin{itemize}
          \item Usa inteiros de 0 a $N-1$.
          \item Semelhante ao slicing de listas.
        \end{itemize}
      \end{block}
      \begin{lstlisting}
df.iloc[0:4, [1, 2]]
      \end{lstlisting}
    \end{column}
  \end{columns}
  \begin{alertblock}{Evite Métodos Antigos}
    O indexador legado \texttt{.ix[]} foi completamente removido nas versões modernas do Pandas.
  \end{alertblock}
\end{frame}

\section{Filtros Booleanos e Sumários}

\begin{frame}[fragile]{Filtros com Máscaras Lógicas}
  \begin{lstlisting}
# Operadores bitwise: & (E), | (OU), ~ (NAO)
# Parenteses sao estritamente OBRIGATORIOS ao redor de cada condicao
pacientes_risco = df[(df['idade'] > 60) & (df['glicose'] >= 126)]
  \end{lstlisting}
  \begin{block}{Inspeção Rápida}
    \begin{itemize}
      \item \texttt{df.head()}: Primeiras 5 linhas.
      \item \texttt{df.info()}: Tipos de dados e contagem de não-nulos.
      \item \texttt{df.describe()}: Média, desvio padrão, quartis, mínimo e máximo.
    \end{itemize}
  \end{block}
\end{frame}

\begin{frame}{Resumo da Aula}
  \begin{itemize}
    \item DataFrames são matrizes heterogêneas com índices nomeados para análise rápida.
    \item Utilize \texttt{.loc} para rótulos e \texttt{.iloc} para posições numéricas.
    \item \textbf{Prática}: Carregamento de CSVs, seleção colunar e filtros booleanos.
  \end{itemize}
\end{frame}

\end{document}
